\subsection{I-22. Two explicit failures of the active-simple-block reduction}
\label{I-22}

\paragraph{Audited claim and scope.}
Singh and Singh (2026), arXiv:2602.13449v1, Theorem~7.4
(printed p.~38), claims the unknown-order bound
$n(wd)^{O(w^2)}$.  Its foundational reductions are false as stated:
the equivalence after Definition~2.1 (p.~4), repeated in
Section~3.1 (p.~6), and the active-block conclusions in
Lemma~3.1 (pp.~6--7), Remark~3.1.4 (p.~9), and Lemma~3.2
(p.~10).  The following counterexamples work over $\mathbb C$,
inside the source's characteristic-zero hypothesis.  They
invalidate those proof steps; they do not disprove the headline
existence of a hitting set with the claimed bound.

\paragraph{Full word algebra loses the prescribed ordered output.}
Let $N=E_{12}$ and $M=E_{21}$, and take the width-two factors
\[
 A_1(x_1)=M+x_1N,\qquad A_2(x_2)=M+x_2N.
\]
Since $N^2=M^2=0$, $NM=E_{11}$, and $MN=E_{22}$,
\[
 A_1(x_1)A_2(x_2)=x_1E_{11}+x_2E_{22}.
\]
The specified ROABP output $e_1^TA_1A_2e_2$ is identically zero.
Nevertheless the algebra generated by the coefficient matrices
is the full $\operatorname{Mat}_2(\mathbb C)$, and
$e_1^TNe_2=1$.  Thus vanishing on the actual ordered products
of one local coefficient per layer does not imply vanishing on
all words in the generated algebra.

This example also separates a nonzero simple-block image of the
matrix product from a nonzero prescribed scalar output.  The
algebra is already simple, so its identity representation is a
simple-block projection and the displayed matrix product survives
it, although the scalar output is zero.  After distinct nonzero
specializations of $x_1,x_2$, the two specialized factors still
generate the full matrix algebra: their difference gives a
nonzero multiple of $N$, and then either factor gives $M$.
Hence moving to a generic multiplicative specialization does not
repair this example.

\paragraph{A nonzero output can live entirely in the radical.}
Take four width-two factors
\[
 A_1=x_1E_{11},\quad A_2=I+x_2N,\quad
 A_3=I+x_3N,\quad A_4=x_4E_{22}.
\]
Direct multiplication gives
\[
 W=A_1A_2A_3A_4=x_1x_4(x_2+x_3)N,\qquad
 f=e_1^TWe_2=x_1x_4(x_2+x_3)\ne0.
 \tag{I.27}
\]
The coefficient algebra is the algebra $\mathcal T_2$ of upper
triangular two-by-two matrices.  Its Jacobson radical is
$J=\mathbb C N$, with $J^2=0$, and its semisimple quotient is
$\mathbb C\oplus\mathbb C$, given by the two diagonal entries.
For completeness, $\mathbb C N$ is a square-zero ideal, hence
lies in the radical; the quotient is a product of fields and
has zero radical, proving equality.  Every homomorphism onto
a simple quotient therefore kills $W$.  No active simple block
of the kind required by the source exists for this nonzero
computation.

All four variables occur, the product length is $4\ge w=2$,
and the example is not a width-one scalar product in disguise:
the coefficient matrix of $x_2+x_3$, with rows and columns
indexed by exponents $0,1$, is
$\left(\begin{smallmatrix}0&1\\1&0\end{smallmatrix}\right)$,
of rank two.  The same obstruction persists after any nonzero
character values for which $x_2+x_3\ne0$; the specialized
algebra remains $\mathcal T_2$ and the product remains a
nonzero radical element.  Constant identity factors can be
inserted to make the product arbitrarily long.

Nilpotency $J^w=0$ controls products containing $w$ radical
factors.  It does not imply that a long product of arbitrary
algebra elements vanishes whenever the final product belongs
to $J$.  Equation~(I.27) gives precisely this distinction.
These two counterexamples are enough to withhold the cited
Theorem~7.4 from the accepted background of this report.
They leave the independently proved I-21 unaffected.
\hfill$\square$
